% ============================================================
\section*{Possibles ajouts dans le modèle}

% -------------------------------------------------------
\subsection*{Terme stérique}

\textit{Pathways to enhance electrochemical CO$_2$ reduction\ldots}

\medskip
Ce terme prend en compte le volume fini des ions (ils ne peuvent pas s'empiler indéfiniment). On ajoute au flux de Nernst--Planck une contribution :

\begin{equation}
    -\mathcal{R} c_i \left(
        \frac{N_a \displaystyle\sum_j a_j^3 \,\nabla c_j}
             {1 - N_a \displaystyle\sum_j a_j^3 \, c_j}
    \right)
\end{equation}

où $\displaystyle\sum_j a_j^3 c_j$ est la fraction volumique occupée par l'espèce $j$, et $1 - \displaystyle\sum_j a_j^3 c_j$ est la fraction volumique libre.

\medskip
Quand $\dfrac{\displaystyle\sum_j a_j^3 \, c_j}{1 - \displaystyle\sum_j a_j^3 \, c_j} \to \infty$, le terme repousse les ions lorsque l'espace libre devient nul.

% -------------------------------------------------------
\subsection*{HER (Hydrogen Evolution Reaction)}

\textit{Pathways to enhance\ldots}

\begin{align}
    2\mathrm{H}^+ + 2e^- &\longrightarrow \mathrm{H}_2 \\
    2\mathrm{H_2O} + 2e^- &\longrightarrow \mathrm{H}_2 + 2\mathrm{OH}^-
\end{align}

\paragraph{Problème :} Le courant total se décompose comme :
\begin{equation}
    j_{\mathrm{tot}} = j_{\mathrm{CO_2RR}} + j_{\mathrm{HER}}
\end{equation}

La HER consomme aussi des H$^+$ et des $e^-$. Elle est négligeable si $j_{\mathrm{HER}} \ll j_{\mathrm{CO_2RR}}$.

% -------------------------------------------------------
\subsection*{Réactions volumiques}

\textit{Pathways to\ldots}

\begin{align}
    \mathrm{CO_2} + \mathrm{H_2O} &\rightleftharpoons \mathrm{H}^+ + \mathrm{HCO_3}^-
        && \rightarrow \text{mauvais pH local} \\
    \mathrm{HCO_3}^- &\rightleftharpoons \mathrm{H}^+ + \mathrm{CO_3}^{2-}
        && \rightarrow \text{mauvaise } c_{\mathrm{CO_2}} \\
    \mathrm{H_2O} &\rightleftharpoons \mathrm{H}^+ + \mathrm{OH}^-
        && \rightarrow \text{électroneutralité incorrecte}
\end{align}

% -------------------------------------------------------
\subsection*{Milieu concentré}

En milieu concentré, on est obligé de passer à l'approche \textbf{Maxwell--Stefan} : on ne peut plus utiliser \textbf{Nernst--Planck} pour calculer $N_i$.

% -------------------------------------------------------
\subsection*{Pas de dépendance en $T$}

\textit{Derivation of an effective thermal\ldots}

On prend $T = \mathrm{cst}$ (sinon effet Soret).

% -------------------------------------------------------
\subsection*{Pas de dépendance en $P$}

On prend $P = \mathrm{cst}$.

% -------------------------------------------------------
\subsection*{Pas de résistance de contact}

On néglige la résistance de contact, ce qui revient à poser $R_c \to 0$ dans :
\begin{equation}
    i_n = \frac{\phi^s - \phi^l}{R_c}
\end{equation}

% -------------------------------------------------------
\subsection*{Formation de bulles}

Le CO$_2$ pourrait buller dans la couche poreuse — cet effet n'est pas pris en compte dans le modèle actuel.